\documentclass[12pt]{amsart}
\usepackage{microtype}
\usepackage{amssymb}
\usepackage{mathtools}
\usepackage{tikz-cd}
\usepackage{mathbbol} % changes \mathbb{} and adds more support
\usepackage{csquotes}
\usepackage{enumitem}
\usepackage{setspace} % For line spacing
\usepackage{mathptmx} % Times New Roman for both text and math
\usepackage[scaled=.92]{helvet} % Helvetica, scaled to match Times New Roman
\usepackage{layout}

% set up 12pt Times New Roman
\renewcommand{\rmdefault}{ptm}
\renewcommand{\sfdefault}{phv}
\renewcommand{\ttdefault}{pcr}

%\usepackage{layout}
\pagestyle{empty}

\setlength{\parindent}{0pt}
\usepackage[letterpaper,
	margin=1.87cm,
	hoffset=0cm,
	voffset=0cm,
	headheight=0cm,
	headsep=0cm,
	marginparwidth=0cm,
	marginparsep=0cm,
	footskip=0cm]{geometry} % Sets paper size and margins

% line spacing
%\onehalfspacing % Alternatively, you can use \singlespacing to try to meet the "no more than six lines per inch" requirement

% bibliography
\usepackage[
	backend=biber,
	style=numeric, % alphabetic, numeric, remove sorting=none
	sorting=nyt,
	backref=false,
	url=false,
	doi=false,
	isbn=false,
	eprint=false]{biblatex}

\setlength{\bibitemsep}{0.0cm}
\renewbibmacro{in:}{}  % don't display "in:" before the journal name
\AtEveryBibitem{\clearfield{pages}}  % don't show page numbers
% environments
\newtheorem{theorem}{Theorem}
\newtheorem*{theorem*}{Theorem}
\newtheorem{proposition}[theorem]{Proposition}
\newtheorem*{proposition*}{Proposition}
\newtheorem{lemma}[theorem]{Lemma}
\newtheorem*{lemma*}{Lemma}
\newtheorem{corollary}[theorem]{Corollary}
\newtheorem*{corollary*}{Corollary}

\theoremstyle{definition}
\newtheorem{definition}[theorem]{Definition}
\newtheorem*{definition*}{Definition}
\newtheorem{remark}[theorem]{Remark}
\newtheorem*{remark*}{Remark}
\newtheorem{example}[theorem]{Example}
\newtheorem*{example*}{Example}
\newtheorem{construction}[theorem]{Construction}
\newtheorem*{construction*}{Construction}
\newtheorem{convention}[theorem]{Convention}
\newtheorem*{convention*}{Convention}
\newtheorem{terminology}[theorem]{Terminology}
\newtheorem*{terminology*}{Terminology}
\newtheorem{notation}[theorem]{Notation}
\newtheorem*{notation*}{Notation}
\newtheorem{question}[theorem]{Question}
\newtheorem*{question*}{Question}

% hyphenation
\hyphenation{co-chain}
\hyphenation{co-chains}
\hyphenation{co-al-ge-bra}
\hyphenation{co-al-ge-bras}
\hyphenation{co-bound-ary}
\hyphenation{co-bound-aries}
\hyphenation{Func-to-rial-i-ty}
\hyphenation{colim-it}
\hyphenation{di-men-sional}

% basics
\DeclareMathOperator{\face}{d}
\DeclareMathOperator{\dege}{s}
\DeclareMathOperator{\bd}{\partial}
\DeclareMathOperator{\sign}{sign}
\newcommand{\ot}{\otimes}
\DeclareMathOperator{\EZ}{EZ}
\DeclareMathOperator{\AW}{AW}

% sets and spaces
\newcommand{\N}{\mathbb{N}}
\newcommand{\Z}{\mathbb{Z}}
\newcommand{\Q}{\mathbb{Q}}
\newcommand{\R}{\mathbb{R}}
\renewcommand{\k}{\Bbbk}
\newcommand{\sym}{\mathbb{S}}
\newcommand{\cyc}{\mathbb{C}}
\newcommand{\Ftwo}{{\mathbb{F}_2}}
\newcommand{\Fp}{{\mathbb{F}_p}}
\newcommand{\Cp}{{\cyc_p}}
\newcommand{\gsimplex}{\mathbb{\Delta}}
\newcommand{\gcube}{\mathbb{I}}

% categories
\newcommand{\Cat}{\mathsf{Cat}}
\newcommand{\Fun}{\mathsf{Fun}}
\newcommand{\Set}{\mathsf{Set}}
\newcommand{\Top}{\mathsf{Top}}
\newcommand{\CW}{\mathsf{CW}}
\newcommand{\Ch}{\mathsf{Ch}}
\newcommand{\simplex}{\triangle}
\newcommand{\sSet}{\mathsf{sSet}}
\newcommand{\cube}{\square}
\newcommand{\cSet}{\mathsf{cSet}}
\newcommand{\Alg}{\mathsf{Alg}}
\newcommand{\coAlg}{\mathsf{coAlg}}
\newcommand{\biAlg}{\mathsf{biAlg}}
\newcommand{\sGrp}{\mathsf{sGrp}}
\newcommand{\Mon}{\mathsf{Mon}}
\newcommand{\symMod}{\mathsf{Mod}_{\sym}}
\newcommand{\symBimod}{\mathsf{biMod}_{\sym}}
\newcommand{\operads}{\mathsf{Oper}}
\newcommand{\props}{\mathsf{Prop}}

% operators
\DeclareMathOperator{\free}{F}
\DeclareMathOperator{\forget}{U}
\DeclareMathOperator{\yoneda}{\mathcal{Y}}
\DeclareMathOperator{\Sing}{Sing}
\newcommand{\loops}{\Omega}
\DeclareMathOperator{\cobar}{\mathbf{\Omega}}
\DeclareMathOperator{\proj}{\pi}
\DeclareMathOperator{\incl}{\iota}
\DeclareMathOperator{\Sq}{Sq}
\DeclareMathOperator{\ind}{ind}

% chains
\DeclareMathOperator{\chains}{N}
\DeclareMathOperator{\cochains}{N^{\vee}}
\DeclareMathOperator{\gchains}{C}

% pair delimiters (mathtools)
\DeclarePairedDelimiter\bars{\lvert}{\rvert}
\DeclarePairedDelimiter\norm{\lVert}{\rVert}
\DeclarePairedDelimiter\angles{\langle}{\rangle}
\DeclarePairedDelimiter\set{\{}{\}}
\DeclarePairedDelimiter\ceil{\lceil}{\rceil}
\DeclarePairedDelimiter\floor{\lfloor}{\rfloor}

% other
\newcommand{\id}{\mathsf{id}}
\renewcommand{\th}{\mathrm{th}}
\newcommand{\op}{\mathrm{op}}
\DeclareMathOperator*{\colim}{colim}
\DeclareMathOperator{\coker}{coker}
\newcommand{\Hom}{\mathrm{Hom}}
\newcommand{\End}{\mathrm{End}}
\newcommand{\coEnd}{\mathrm{coEnd}}
\newcommand{\xla}[1]{\xleftarrow{#1}}
\newcommand{\xra}[1]{\xrightarrow{#1}}
\newcommand{\defeq}{\stackrel{\mathrm{def}}{=}}

% comments
\newcommand{\anibal}[1]{\noindent\textcolor{blue}{\underline{Anibal}: #1}}

% pdf
\newcommand{\pdfEinfty}{\texorpdfstring{${E_\infty}$}{E-infty}}

% mathrm
\newcommand{\rA}{\mathrm{A}}
\newcommand{\rB}{\mathrm{B}}
\newcommand{\rC}{\mathrm{C}}
\newcommand{\rD}{\mathrm{D}}
\newcommand{\rE}{\mathrm{E}}
\newcommand{\rF}{\mathrm{F}}
\newcommand{\rG}{\mathrm{G}}
\newcommand{\rH}{\mathrm{H}}
\newcommand{\rI}{\mathrm{I}}
\newcommand{\rJ}{\mathrm{J}}
\newcommand{\rK}{\mathrm{K}}
\newcommand{\rL}{\mathrm{L}}
\newcommand{\rM}{\mathrm{M}}
\newcommand{\rN}{\mathrm{N}}
\newcommand{\rO}{\mathrm{O}}
\newcommand{\rP}{\mathrm{P}}
\newcommand{\rQ}{\mathrm{Q}}
\newcommand{\rR}{\mathrm{R}}
\newcommand{\rS}{\mathrm{S}}
\newcommand{\rT}{\mathrm{T}}
\newcommand{\rU}{\mathrm{U}}
\newcommand{\rV}{\mathrm{V}}
\newcommand{\rW}{\mathrm{W}}
\newcommand{\rX}{\mathrm{X}}
\newcommand{\rY}{\mathrm{Y}}
\newcommand{\rZ}{\mathrm{Z}}
% mathcal
\newcommand{\cA}{\mathcal{A}}
\newcommand{\cB}{\mathcal{B}}
\newcommand{\cC}{\mathcal{C}}
\newcommand{\cD}{\mathcal{D}}
\newcommand{\cE}{\mathcal{E}}
\newcommand{\cF}{\mathcal{F}}
\newcommand{\cG}{\mathcal{G}}
\newcommand{\cH}{\mathcal{H}}
\newcommand{\cI}{\mathcal{I}}
\newcommand{\cJ}{\mathcal{J}}
\newcommand{\cK}{\mathcal{K}}
\newcommand{\cL}{\mathcal{L}}
\newcommand{\cM}{\mathcal{M}}
\newcommand{\cN}{\mathcal{N}}
\newcommand{\cO}{\mathcal{O}}
\newcommand{\cP}{\mathcal{P}}
\newcommand{\cQ}{\mathcal{Q}}
\newcommand{\cR}{\mathcal{R}}
\newcommand{\cS}{\mathcal{S}}
\newcommand{\cT}{\mathcal{T}}
\newcommand{\cU}{\mathcal{U}}
\newcommand{\cV}{\mathcal{V}}
\newcommand{\cW}{\mathcal{W}}
\newcommand{\cX}{\mathcal{X}}
\newcommand{\cY}{\mathcal{Y}}
\newcommand{\cZ}{\mathcal{Z}}
% mathsf
\newcommand{\sA}{\mathsf{A}}
\newcommand{\sB}{\mathsf{B}}
\newcommand{\sC}{\mathsf{C}}
\newcommand{\sD}{\mathsf{D}}
\newcommand{\sE}{\mathsf{E}}
\newcommand{\sF}{\mathsf{F}}
\newcommand{\sG}{\mathsf{G}}
\newcommand{\sH}{\mathsf{H}}
\newcommand{\sI}{\mathsf{I}}
\newcommand{\sJ}{\mathsf{J}}
\newcommand{\sK}{\mathsf{K}}
\newcommand{\sL}{\mathsf{L}}
\newcommand{\sM}{\mathsf{M}}
\newcommand{\sN}{\mathsf{N}}
\newcommand{\sO}{\mathsf{O}}
\newcommand{\sP}{\mathsf{P}}
\newcommand{\sQ}{\mathsf{Q}}
\newcommand{\sR}{\mathsf{R}}
\newcommand{\sS}{\mathsf{S}}
\newcommand{\sT}{\mathsf{T}}
\newcommand{\sU}{\mathsf{U}}
\newcommand{\sV}{\mathsf{V}}
\newcommand{\sW}{\mathsf{W}}
\newcommand{\sX}{\mathsf{X}}
\newcommand{\sY}{\mathsf{Y}}
\newcommand{\sZ}{\mathsf{Z}}
% mathbb
\newcommand{\bA}{\mathbb{A}}
\newcommand{\bB}{\mathbb{B}}
\newcommand{\bC}{\mathbb{C}}
\newcommand{\bD}{\mathbb{D}}
\newcommand{\bE}{\mathbb{E}}
\newcommand{\bF}{\mathbb{F}}
\newcommand{\bG}{\mathbb{G}}
\newcommand{\bH}{\mathbb{H}}
\newcommand{\bI}{\mathbb{I}}
\newcommand{\bJ}{\mathbb{J}}
\newcommand{\bK}{\mathbb{K}}
\newcommand{\bL}{\mathbb{L}}
\newcommand{\bM}{\mathbb{M}}
\newcommand{\bN}{\mathbb{N}}
\newcommand{\bO}{\mathbb{O}}
\newcommand{\bP}{\mathbb{P}}
\newcommand{\bQ}{\mathbb{Q}}
\newcommand{\bR}{\mathbb{R}}
\newcommand{\bS}{\mathbb{S}}
\newcommand{\bT}{\mathbb{T}}
\newcommand{\bU}{\mathbb{U}}
\newcommand{\bV}{\mathbb{V}}
\newcommand{\bW}{\mathbb{W}}
\newcommand{\bX}{\mathbb{X}}
\newcommand{\bY}{\mathbb{Y}}
\newcommand{\bZ}{\mathbb{Z}}
% mathfrak
\newcommand{\fA}{\mathfrak{A}}
\newcommand{\fB}{\mathfrak{B}}
\newcommand{\fC}{\mathfrak{C}}
\newcommand{\fD}{\mathfrak{D}}
\newcommand{\fE}{\mathfrak{E}}
\newcommand{\fF}{\mathfrak{F}}
\newcommand{\fG}{\mathfrak{G}}
\newcommand{\fH}{\mathfrak{H}}
\newcommand{\fI}{\mathfrak{I}}
\newcommand{\fJ}{\mathfrak{J}}
\newcommand{\fK}{\mathfrak{K}}
\newcommand{\fL}{\mathfrak{L}}
\newcommand{\fM}{\mathfrak{M}}
\newcommand{\fN}{\mathfrak{N}}
\newcommand{\fO}{\mathfrak{O}}
\newcommand{\fP}{\mathfrak{P}}
\newcommand{\fQ}{\mathfrak{Q}}
\newcommand{\fR}{\mathfrak{R}}
\newcommand{\fS}{\mathfrak{S}}
\newcommand{\fT}{\mathfrak{T}}
\newcommand{\fU}{\mathfrak{U}}
\newcommand{\fV}{\mathfrak{V}}
\newcommand{\fW}{\mathfrak{W}}
\newcommand{\fX}{\mathfrak{X}}
\newcommand{\fY}{\mathfrak{Y}}
\newcommand{\fZ}{\mathfrak{Z}}

\addbibresource{../usualpapers.bib}
\addbibresource{../bibliography.bib}
% !TEX root = ../proposal.tex

\newcommand{\coho}{\rH^\bullet}
\renewcommand{\cochains}{\rC^\bullet}
\newcommand{\KH}{\mathrm{KH}_\bullet}
\newcommand{\giottoTDA}{\raisebox{-.5pt}{\includegraphics[scale=.08]{aux/giotto}}}
\DeclareMathOperator{\ad}{ad}
\DeclareMathOperator{\barconst}{\mathbf{B}}
\newcommand{\desus}[1]{s^{\mbox{\scriptsize -#1}}}
\DeclareMathSymbol{\antishrik}{\mathord}{operators}{"3C}
\newcommand{\shrik}{!}
\newcommand{\cdga}{\mathsf{cdga}}
\newcommand{\cdgl}{\mathsf{cdgl}}
\newcommand{\com}{\mathcal{C}om}
\newcommand{\lie}{\mathcal{L}ie}
\newcommand{\Der}{\mathrm{Der}}
\newcommand{\fullline}{\par\noindent\rule{\textwidth}{0.4pt}}


\usepackage{csquotes}
\usepackage{enumitem}
\usepackage[letterpaper, margin=1.87cm]{geometry} % Sets paper size and margins
\usepackage{setspace} % For line spacing
\usepackage{mathptmx} % Times New Roman for both text and math
\usepackage[scaled=.92]{helvet} % Helvetica, scaled to match Times New Roman

% set up 12pt Times New Roman
\renewcommand{\rmdefault}{ptm}
\renewcommand{\sfdefault}{phv}
\renewcommand{\ttdefault}{pcr}

% line spacing
%\onehalfspacing % Alternatively, you can use \singlespacing to try to meet the "no more than six lines per inch" requirement


%%%%%%%%%%%%%%%%%%%%%%%%%%%%%%%%%%%%%5

\title[Effective algebro-homotopical constructions]{Effective algebro-homotopical constructions \\ with applications to data science and quantum field theory}

\author{Anibal~M.~Medina-Mardones}
%\address{A.M-M., Max Planck Institute for Mathematics \and University of Notre Dame}
%\email{\href{mailto:ammedmar@mpim-bonn.mpg.de}{ammedmar@mpim-bonn.mpg.de}}

\begin{document}
	\vspace*{-0pt}
	\maketitle
	\thispagestyle{empty}
	% !TEX root = ../proposal.tex

\section*{Introduction} \label{s:introduction}

There is a tense trade-off in algebraic topology having roots reaching back to the beginning of its modern form.
This tension can be illustrated with the concept of cohomology.
The first approaches, dating back to Poincar\'e, are based on the subdivision of a space into simple pieces.
These elementary shapes are made to generate a free graded module and their spatial relations define the differential used to compute cohomology.
This definition makes certain geometric properties of cohomology, for example excision, fairly clear.
Yet, it is not easy to show that a continuous map of spaces induces a map between their associated cohomologies.
The functoriality just alluded to is trivial when defining cohomology in terms of homotopy classes of maps to Eilenberg--MacLane spaces, but the passage to the homotopy category erases geometric and combinatorial information and the resulting definition is not well suited for concretely presented spaces.
The tension this example illustrates manifests itself in many other important contexts, and the trade-off between concreteness and functoriality is as central today as it was almost a century ago.

\section*{Objectives}

My long-term vision is to ease this tension by developing concrete constructions of homotopical invariants currently defined only indirectly; and, by doing so, extending the reach of algebraic topology into fields beyond pure mathematics, including Data Science and Quantum Field Theory.
We will organize the presentation of this program around four interrelated deliverables.

\smallskip
\quad (1) Develop software to compute secondary operations on triangulated spaces.\par
\quad (2) Compute all Steenrod operations on Khovanov homology for knots with up to 14 crossings.\par
\quad (3) Compute the Steenrod barcodes of densely sampled molecular conformation spaces.\par
\quad (4) Relate the sate sum and functorial classifications of invertible topological phases of matter.\par
	% !TEX root = ../proposal.tex

\section*{(1) Effective Algebro-Homotopical Constructions and Their Implementations \\ (PhD1, MSc1, MSc2)}

A foundational goal of the systematic study of topological spaces up homotopy equivalence is the development of algebraic invariants, such as cohomology.
As mentioned in the introduction, Poincar\'e's definition of cohomology $\rH^\bullet$ relies on the cochains $\cochains$ of a cellular space.
A ground-breaking result in modern homotopy theory, due to Mandell \cite{mandell2006homotopy_type}, is that, under certain finiteness assumptions, all the homotopical information of a cellular space is captured by a so-called $E_\infty$-algebra structure on $\cochains$.
Aided by this language, the first goal of our program can be more precisely stated as follows:

\smallskip\noindent(1a) Effectively describing progressively richer layers of homotopical data using the $E_\infty$-algebra on $\cochains$.\par
\noindent(1b) Implementing these constructions into computer algebra systems (CAS) such as \texttt{ComCH}, \texttt{SAGE}, and \texttt{Maple}.

\smallskip\noindent These developments will render the invariants within these layers accessible through computer-based computations, which is indispensable for their use in data science (Section 3), and, additionally, they will play an essential role in providing local descriptions of topological action functionals on triangulated space-times (Section 4).

\medskip\noindent\textbf{Earlier work}.
A fundamental step in the direction of our goal is the description of a ring structure on $\rH^\bullet$ using the Alexander--Whitney product on simplicial cochains $\cochains$ \cite{alexander1936ring, whitney1938products}.
Let us illustrate with an example the additional information captured by this ring structure.
The cohomology groups of $\mathbb{C} P^2$ and $S^2 \vee S^4$, the union of a $2$- and a $4$-sphere over a point, are isomorphic despite these spaces not being homotopy equivalent.
By choosing triangulations for these spaces, the Alexander--Whitney product can be used to show that the ring structure in cohomology distinguishes them.
An algorithm for this type of computation is available for example on \texttt{SAGE}.
Considering the suspension of these spaces leads to two non-homotopic spaces $\Sigma(\mathbb{C} P^2)$ and $\Sigma(S^2 \vee S^4)$ whose cohomologies are isomorphic as graded rings.
Distinguishing these effectively lead us to a novel step in the direction of our goal.
Let us consider a prime $p$.
Steenrod made $\rH^\bullet$ it into a module over his algebra $\cA_p$ \cite{steenrod1962cohomology}.
This structure is explicitly induced from the $E_\infty$-algebra structure on simplicial cochains by means of the cup-$(p,i)$ products.
These generalize to odd primes the cup-$i$ products of Steenrod, and were introduced and implemented, using effective versions of May's operadic methods \cite{may1970general}, in \cite{medina2021may_st,medina2021comch}.
The action of $\cA_2$ suffices to distinguish $\Sigma(\mathbb{C} P^2)$ and $\Sigma(S^2 \vee S^4)$, but and odd prime will be needed for similar comparisons using the suspension of quaternionic projective spaces, for example.

%Steenrod operations in the mod $p$ cohomology of spaces, which together with the Bockstein homomorphism provide a complete account of the mod $p$ cohomology functor.

\medskip\noindent\textbf{{\sc Relations}}.
The next layer of homotopical structure after Steenrod operations comes from the relations these operations satisfy.
By group homology computations, or more abstractly Mandell's Theorem, we known that there must be some structure on the $E_\infty$-algebra on $\cochains$ inducing these so-called \textit{Cartan} and \textit{Adem relations}.
%We referred to these as Cartan and Adem coboundaries.
Recently, these structures have been effectively described over any prime for Cartan's, and over the even prime for Adem's \cite{medina2020cartan,medina2023oddcartan,medina2021adem}.

\medskip\noindent\textbf{Goals}.
My team will culminate this research direction and consolidate it into useful software.
More specifically, we will complete the following tasks:

\smallskip\noindent(MSc1) Implementing the existing Cartan and Adem structures over the even prime.\par
\noindent(MSc2) Implementing the existing Cartan structure over odd primes.\par
\noindent(PhD1) Discovering and implementing the Adem structure over odd primes.

\smallskip\noindent As we will detail in Section 4, these structures have already proven important in the classification of topological phases of matter.
Additionally, making them accessible through computer computations represents the initial stride toward integrating secondary operations into persistent cohomology, as discussed in Section 3.

\smallskip\noindent\textbf{Methodology}.
We have already established a strategic partnership with the Ontario Research Centre for Computer Algebra (ORCCA) and all students in the my program will actively participate in our weekly meetings.

(MSc1) will review the existing implementation of the mod 2 Cartan and Adem coboundary constructor prototyped in \texttt{Python}.
They will then develop a compatible framework for this routine to be compatible with the large scale project $\texttt{SAGE}$.

(MSc2) will use this infrastructure to design and implement an odd prime Cartan coboundary constructor directly into \texttt{SAGE}.

(PhD1) will use the new models of the $E_\infty$-operad \cite{medina2020prop1} to prove the odd Adem relation effectively.
This will require a solid grasp of group homology, which will be provided in the form of course taught by the PI for the department at large.
Following this significant achivement, (PhD1) will use the code base developed by (MSc1) and (MSc2) to design and implement into \texttt{SAGE} a constructor of odd prime Adem coboundaries.

\medskip\noindent\textbf{{\sc Khovanov homology} (PhD2)}
$\mathrm{KH_\bullet}$ is a powerful algebraic invariant of knots and links which refines the Jones polynomial \cite{khovanov2000khovanov}.
It is effectively computable from a chain complex associated with a knot diagram, and several implementations of algorithms for this task exist.
A ground-breaking result of two teams (\cite{lipshitz2014khovanov,kriz2016khovanov}) states that this invariant can be obtained from a cellular spectrum; which implies, indirectly, that $\mathrm{KH_\bullet}$ has an action of the Steenrod algebra $\cA_p$ for any prime $p$.
My team will focus on the design and implementation of algorithms computing this finer invariant effectively.

\medskip\noindent\textbf{Earlier work}.
Adapting the formulas of \cite{medina2023fast_sq}, Cantero-Mor\'an constructed structure at the chain level inducing the action of $\cA_2$ on $\KH$ \cite{cantero-moran2020khovanov}.
His effective methods recover those of Lipshitz and Sakar, which deal with the computation of $\Sq^1$ and $\Sq^2$ \cite{lipshitz2014steenrod}.
These two operations were implemented by Seed \cite{seed2012khovanov} and, last year, all squares were implemented into \texttt{SAGE} by Milstein using Cantero-Morán's structure \cite{milstein2022khovanov}.
Put together, these developments have effectively provided a complete treatment of all primary operations on $\KH$ at the even prime.

\medskip\noindent\textbf{Goals}.
Provide a complete treatment of primary operations on $\KH$ at all primes.
Explicitly, we will achieve the following goals:

\smallskip\noindent(PhD2a)
Effectively describing structure inducing the action on $\KH$ of $\cA_p$ for any prime p.\par
\smallskip\noindent(PhD2b)
Designing and implementing algorithms for the computation of all Steenrod operations on $\KH$.\par
\smallskip\noindent(PhD2c)
Computing these invariants for $p=3,5,7$ and knots with up to 14 crossing.

%\smallskip\noindent This software's primary impact will be in the field of low-dimensional topology, enabling a more precise exploration of intricate properties related to knots and links.

\smallskip\noindent\textbf{Methodology}.
For this project, (PhD2) will work closely with both the PI and Cantero-Morán, who will serve as a secondary advisor.
The first step of this project, to be accomplished as joint work by Cantero-Mor\'an, (PhD2), and the PI, will generalize the PI's cup-$i$ formulas \cite{medina2023oddcartan} to the odd prime case.
Preliminary work as shown that this task is complex and interesting on its own right, being related to Tate cohomology of the symmetric groups.
This preparatory step is needed since the existing cup-$(p,i)$ formulas, introduced by the PI in \cite{medina2021may_st}, are not adequate for the adaptation to the Burnside category done by Cantero-Mor\'an in the even prime case \cite{cantero-moran2020khovanov}.
After this step is completed, (PhD2) will generalize said adaptation for odd primes and design an algorithm for the computation of these invariants.
Finally, (PhD2) will develop, in consultation with other members of ORCCA, a high-performance tool for the computation of these invariants from diagram presentations, and produce a catalogue of Steenrod operations for all knots up to 14 crossings.

	% !TEX root = ../proposal.tex

\section*{(2) Refinements and Novel Applications of Persistent Cohomology \\ (PhD3, MSc3, MSc4)}

Persistent (co)homology is a powerful technique used in the field of topological data analysis.
It serves as a valuable tool for analyzing and extracting topological features and structures from data, especially in the context of point clouds or other high-dimensional data representations.
The core steps i the persistent (co)homology pipeline are the following:
i) Through various standard techniques, construct from the given data triangulated spaces with a filtration
\[
X_0 \subseteq X_1 \subseteq X_2 \subseteq \dots \subseteq X_N \,.
\]
ii) Apply the (co)homology functor to obtain a diagram
\[
\rH^\bullet X_0 \leftarrow \rH^\bullet X_1 \leftarrow \dots \leftarrow \rH^\bullet X_N \,.
\]
iii) Summarize the resulting sequence of homologies into a computationally efficient invariant known as the ``barcode" of the filtration.
A key feature of persistent (co)homology is the stability property, which asserts that, with appropriate metrics, the passage to barcodes maintains bounded distortion.

Numerous high-performance implementations of algorithms for computing barcodes are available, such as \texttt{giotto-tda} \cite{medina2021giotto}, \texttt{ripser} \cite{bauer2021ripser}, and \texttt{gudhi} \cite{gudhi}.
This is one of the main reasons explaining the significant impact this technique is having across scientific domains.
Consult \cite{carlsson2021topological} for a fairly recent survey.
However, it is important to recognize that the classical barcode, despite its utility, is not without substantial limitations, as it does not encode the way bars interact with each other.
These limitations extend from those associated with (non-persistent) cohomology discussed in the previous section.

The second goal of this research program is to develop finer computable invariants for persistent (co)homology, provide high-performance implementations for the computation of new and existing invariants, and identify domains of applicability where this techniques will proved domain specific insights.
These software tools will enable the analysis of significantly larger datasets, offering ample examples to interpret these innovative topological signature of data within domain-specific contexts, such as molecular chemistry and neuroscience.

\medskip\noindent\textbf{Earlier work}.
Using the novel algorithm for the calculation of Steenrod squares introduced in \cite{medina2023fast_sq}, we effectively integrated finer information encoded by Steenrod operations, specifically, the action of $\cA_2$ on mod 2 cohomology, into the persistence (co)homology pipeline \cite{medina2022per_st}.
The resulting invariant takes the form of a sequence of barcodes for $k \geq 0$, with $k = 0$ recovering the traditional mod 2 barcode. For $k > 0$, these barcodes are also stable and encode generalizations of self-intersection of cycles on a closed manifold. Importantly, they are computable.
In fact, in collaboration with members of the \giottoTDA\ team, we developed \texttt{steenroder}, a tool for the computation of Steenrod barcodes.
This tool played a crucial role in detecting the presence of Steenrod barcodes in real-world data, as detailed in \cite{medina2022per_st}, specifically within the conformation space of the cyclo-octane molecule.
However, it is worth noting that the performance of this tool is not yet optimal.
In the following section, we will describe how to improve its efficiency.

In another significant conceptual advance towards enhancing persistent (co)homology with finer computable invariants, as discussed in \cite{memoli2022cup_length}, the authors introduced the notion of the cup-length invariant. This innovation incorporates a portion of the ring structure on cohomology into the persistent homology pipeline.
More precisely, the cup-length of a graded ring is defined as the largest non-negative integer $\ell$ for which there exist positive-dimensional homogeneous elements $a_1, \dots, a_\ell$ within the ring such that their product, $a_1 \dots a_\ell$, is nonzero. The persistent version of this invariant is defined as a function over the filtration parameter.
Regrettably, it is worth noting that as of now, there is no existing implementation of their algorithm.

\medskip\noindent\textsc{High-performance \texttt{steenroder} and \texttt{cuplengther}}
There are three steps in the computation of Steenrod barcodes:
i) The first step involves computing the traditional barcode. This is accomplished through a matrix reduction routine that has been extensively optimized through the collaborative efforts of multiple researchers.
For a cross-software performance comparison one can consult \cite{otter2017roadmap}.
ii)~The second step utilizes a representative cocycle $\alpha$ for each bar to construct a representative cocycle $\beta$ for $\Sq^k[\alpha]$.
This is done using the algorithm introduced in \cite{medina2023fast_sq}.
iii)~Finally, the third step involves another matrix reduction procedure that examines the linear independency of the resulting Steenrod cocycles.
The second step, the current bottleneck of \texttt{steenroder}, is quadratic on the length of the cocycle representative $\alpha$.
The delays are caused by overly representatives with large support, and the lack of parallelization of the routine across cocycles.

Additionally, advances on cocycle replacement will allow for the efficient implementation of an algorithm computing the persistent cup-length invariant.

\medskip\noindent\textbf{Goals}.
Improve the performance of \texttt{steenroder}, implement \texttt{cuplengther}, and test them on real-world data.
Explicitly, my team will complete the following tasks:

\smallskip\noindent(MSC3a)
Parallelizing the construction of Steenrod cocycles across cocycle representatives.

\smallskip\noindent(PhD3a)
Replacing naively obtained cocycles with alternatives having smaller support.

\smallskip\noindent(PhD3b) Implementing \texttt{cuplengther}, a tool for the computation of persistent cup-length.

\smallskip\noindent(PhD3c)
Obtaining a catalogue of topological signatures of densely sampled molecular conformation spaces.

\smallskip\noindent\textbf{Methodology}.
The software development components of this project will be conducted in collaboration with ORCCA for overall support.
Additionally, for the specific application in chemistry, we will partner with the Theory and Computation research group at Western.

	% !TEX root = ../proposal.tex

\smallskip
{\centering \underline{\textsc{Build a bridge between state sum and functorial definitions of TQFTs}}\par}
{\centering Objective (4) \hspace*{2cm} --- \hspace*{2cm} (MSc3, PostDoc)\par}

\medskip Let us return to the trade-off between constructiveness and functoriality discussed in the introduction, now in the context of quantum physics.

\smallskip\noindent\textbf{Background}.
The significance of topology in condensed matter physics is exemplified by the 2016 Nobel Prize in Physics, awarded for ``theoretical discoveries of topological phase transitions and topological phases of matter."
There are primarily two approaches to defining these TQFT's.
On one hand, they are presented using \textit{lattice models}, which can extended to more general space-times using manifold triangulations.
An important class of these are the so-called \textit{symmetry-protected invertible topological phases}.
Using cochain level structures including cup-$i$ products, and Adem and Cartan coboundaries, these have been classified in low dimensions by Kitaev, Gaiotto--Kapustin, and other physicists \cite{kitaev2009periodic, kapustin2015cobordism, barkeshli2022classification}.
Their examples support the ansatz that such classification is controlled by generalized cohomology theories \cite{adams1995stable}.

\quad On the other hand, the functorial approach to TQFTs, pioneered by Atiyah and Segal \cite{atiyah1988tqft}, organizes the discussion around a higher categorical version of functors from the category of manifolds and bordisms to the category of complex vector spaces \cite{baez1995higher}.
In the celebrated work of Hopkins and Freed \cite{freed2021reflection}, this functorial viewpoint was used to relate the classification of invertible topological phases to stable homotopy theory \cite{freed2021reflection}, in line with what physicists expected, but the methods are wildly different, and the way the higher categorical theory emerges from the effective methods on discrete space-times is a key question in the field.
We quote from a lecture given by Hopkins in Regensburg \cite{hopkins2023higher}:
``\textit{This story is far from being understood, and there is a big arc of research that goes into trying to really derive this emerging structure from lattice models, and this is very interesting and very difficult area right now."}

\smallskip\noindent\textbf{Strategic partnership}.
This research objective will be accomplished in collaboration with researchers from the Perimeter Institute, including L. M\"uller and A. Turzillo.

\smallskip\noindent\textbf{Key idea}.
What sets our approach apart from others—and serves as a unifying thread—is our use of Street's free $n$-category on the $n$-simplex \cite{street1987orientals}, commonly referred to as the $n^\th$ oriental $\cO_n$.
The aptness of this object for the stated goal arises from its grounding on both the discrete and higher categorical settings.
% grounded in
Using these canonical categories, one can define the $n$-category generated by a triangulated space $X$ as $\cO(X) = \colim_{\gsimplex\, \downarrow X} \cO_n$.
The connection between Steenrod's cup-$i$ products, used in the lattice formulas, and Street's orientals was established in \cite{medina2020globular}, where the author demonstrated that the intricate structure of $\mathcal{O}_n$ can be naturally derived from the cup-$i$ products of Steenrod.

\smallskip\noindent\textbf{Methodology}.
When $X$ is a triangulation of an $n$-manifold $M$, we will enhance the above colimit to reflect tangential structure on $M$, with orientations and framings being the key examples.
We will introduce the notion of \textit{functorial state sum data} as a cone under the tangentially enhanced Street diagram, whose apex is a sufficiently dualizable $n$-category.
We will then obtain invariants like the partition function using the colimit description of $\cO(X)$.
Preliminary investigations have shown that this method recovers known state sum descriptions of low dimensional TQFTs, including Turaev--Viro's, and, given the systematic nature of our approach, it will enable us to define new theories, including a non-invertible version of Dijkgraaf--Witten theory.

\quad The techniques developed to accomplish this research objective will also have an impact on theoretical mathematics.
For example, my student Aaron Huntley (\underline{MSc6}) will prove the following conjecture emanating from our preliminary work towards this objective: An $A_\infty$-algebra, i.e., a representation of the Stasheff operad in the category of chain complexes $\Ch$, is the exact same data as a cone over the traditional (i.e. non-tangentially enhanced) Street diagram with apex the desuspension $\rB\Ch$.

\smallskip\noindent\textbf{Significance}.
Achieving a comprehensive theory of functorial state sum data, along with a deeper understanding of its connection to cochain-level structures, will enable us to reconcile the functorial and state-sum approaches to classifying invertible topological phases.
This will not only allow for the definition of new TQFTs, but also elucidate the divergent paths taken by Gaiotto--Kapustin and Freed--Hopkins in their respective approaches.
	\newpage
	\sloppy
	\printbibliography
\end{document}

Proposal

Five pages max.

Addressing the points below, describe the proposed research to be supported. Images and graphics are included in the page limit.

Recent progress

Describe your recent progress in research activities related to the proposal; for returning grantees, describe as well the progress attributable to your previous Discovery Grant.

Objectives

Define the short- and long-term objectives of your research program. Note that a research program should have a long-term vision that expands beyond the five years of the Discovery Grant. A single, short-term project or collection of projects does not constitute a research program.

Literature review

Discuss the literature pertinent to the proposal, placing the proposed research in the context of the state of the art.

Methodology

Describe the methods and proposed approach, providing sufficient details to allow the reviewers to assess the feasibility of the research activities.

Considering equity, diversity and inclusion (EDI) in the research process promotes research excellence by making research outcomes more ethically sound, rigorous, reproducible, and useful. It is important to consider EDI through each stage of the research process including, but not limited to, the research questions, design, methodology, analysis, interpretation and dissemination of results, and integrate these considerations where relevant. Consult Equity, diversity and inclusion considerations at each stage of the research process for more information.

Impact

Explain the anticipated significance of the work.

Note:

If the information provided is insufficient, NSERC reserves the right to take appropriate action, such as not soliciting reports from external reviewers or withdrawing the application from the competition.
If relevant to your research, consult NSERC’s Guidelines for the preparation and review of applications in interdisciplinary research and/or Guidelines for the preparation and review of applications in engineering and the applied sciences.
If NSERC determines that the subject matter is outside of its mandate at any time during the review cycle, the application will be rejected. See This link will take you to another Web site Selecting the appropriate federal granting agency for more information.
